A bijective linear mapping $\varphi: E \to F$ is called a linear isomorphism and is denoted by $\varphi: E \xrightarrow{\cong} F$.

Given a linear isomorphism $\varphi: E \to F$, the inverse set mapping $\varphi^{-1}: E \leftarrow F$ is also a linear mapping and hence a linear isomorphism. It is called the inverse isomorphism of $\varphi$.

Two vector spaces $E$ and $F$ are called isomorphic if there exists a linear isomorphism from $E$ onto $F$.
